% =============================================================================
%  bda_preamble.tex
%  Shared LaTeX preamble for Business Data Analytics course units.
%  -----------------------------------------------------------------------------
%  Loaded by every Unit's lecture.Rnw via
%    % =============================================================================
%  bda_preamble.tex
%  Shared LaTeX preamble for Business Data Analytics course units.
%  -----------------------------------------------------------------------------
%  Loaded by every Unit's lecture.Rnw via
%    % =============================================================================
%  bda_preamble.tex
%  Shared LaTeX preamble for Business Data Analytics course units.
%  -----------------------------------------------------------------------------
%  Loaded by every Unit's lecture.Rnw via
%    % =============================================================================
%  bda_preamble.tex
%  Shared LaTeX preamble for Business Data Analytics course units.
%  -----------------------------------------------------------------------------
%  Loaded by every Unit's lecture.Rnw via
%    \input{../../style_references/bda_preamble.tex}
%  immediately after \documentclass[17pt]{extarticle}.
%
%  Unit-specific configuration that must remain in lecture.Rnw, after this
%  \input but before \begin{document}:
%    \fancyhead[L]{\small\color{hdrBG}\textbf{Unit\ N <chapter title>}}
%    \setcounter{theorem}{0}\setcounter{example}{0}
%    \renewcommand\thedefinition{N.\arabic{theorem}}
%    \renewcommand\thetheorem  {N.\arabic{theorem}}
%    \renewcommand\theexample  {N.\arabic{example}}
%    \renewcommand\theexercise {N.\arabic{exercise}}
%    \renewcommand\thefigure   {N.\arabic{figure}}
%    \title{\color{hdrBG}\bfseries\Huge <chapter title>}\author{}\date{}
% =============================================================================

\usepackage[
  paperwidth=257mm,
  paperheight=144.5mm,
  top=9mm,bottom=9mm,left=6mm,right=6mm,
  headheight=7mm,headsep=2mm,footskip=19pt
]{geometry}

\usepackage{amsmath,amsthm,amssymb,amsfonts}
\usepackage{xstring}  % \IfDecimal used by ctab for numeric-vs-font dispatch
\usepackage{booktabs,array,longtable,tabularx}
\usepackage[shortlabels]{enumitem}
\setlist[itemize,1]{leftmargin=1.5em, labelsep=0.35em}
\usepackage{mathtools}
\usepackage{xcolor}
\usepackage{fancyhdr}
\usepackage{tikz}
\usetikzlibrary{arrows.meta,calc,decorations.markings}
\usepackage[most]{tcolorbox}
\tcbuselibrary{theorems,skins,breakable}
\usepackage[round,authoryear]{natbib}  % loaded before hyperref so \citep is hyperlinked
\usepackage{hyperref}
\hypersetup{colorlinks=true,linkcolor=blue!60!black,urlcolor=blue!70!black,citecolor=blue!60!black}

% Source code highlighting (requires pygmentize; compile with -shell-escape)
\usepackage[cache=false]{minted}
\setminted{fontsize=\small, frame=single, framesep=2mm, rulecolor=gray!40,
  bgcolor=gray!5, breaklines=false, tabsize=2,
  autogobble=true,
  linenos=true,
  numbersep=5pt,
  xleftmargin=2.5em}

% Line-number style: gray, slightly smaller, so it does not compete with the code.
\renewcommand{\theFancyVerbLine}{%
  \textcolor{gray!60}{\scriptsize\arabic{FancyVerbLine}}}

% xeCJK is loaded for safety even though all course content is in English;
% it lets occasional CJK characters (e.g. a student name in a sample dataset)
% render without breaking the build.
\usepackage{xeCJK}
\setCJKmainfont{Noto Serif CJK TC}
\setCJKsansfont{Noto Sans CJK TC}
\setCJKmonofont{Noto Sans Mono CJK TC}

% Course-wide color palette.
\definecolor{hdrBG}{HTML}{1E3A8A}
\definecolor{accent}{HTML}{0EA5E9}
\definecolor{defBG}{HTML}{EFF6FF}\definecolor{defFR}{HTML}{2563EB}
\definecolor{thmBG}{HTML}{FEF3C7}\definecolor{thmFR}{HTML}{B45309}
\definecolor{exaBG}{HTML}{FEF2F2}\definecolor{exaFR}{HTML}{B91C1C}
\definecolor{proofBG}{HTML}{F8FAFC}\definecolor{proofFR}{HTML}{64748B}
\definecolor{exBG}{HTML}{FFF7ED}\definecolor{exFR}{HTML}{C2410C}
\definecolor{solBG}{HTML}{F0FDF4}\definecolor{solFR}{HTML}{15803D}
\definecolor{rmkBG}{HTML}{F1F5F9}\definecolor{rmkFR}{HTML}{94A3B8}

% Header / footer skeleton. \fancyhead[L] is set per-unit in lecture.Rnw.
\pagestyle{fancy}\fancyhf{}
\renewcommand{\headrulewidth}{0.4pt}
\renewcommand{\headrule}{\color{hdrBG!30}\hrule width\textwidth height0.4pt}
\fancyhead[R]{\small\color{hdrBG}\leftmark}
\fancyfoot[R]{\small\color{gray}\thepage}
\renewcommand{\sectionmark}[1]{\markboth{\S\thesection~#1}{}}

% Theorem-like environments. All use \theoremstyle{definition} so the body
% renders upright (no italic / slanted text in theorem statements, proofs, or
% solutions).
\theoremstyle{definition}
\newtheorem{theorem}{Theorem}
\newtheorem{lemma}[theorem]{Lemma}
\newtheorem{definition}[theorem]{Definition}
\newtheorem{example}{Example}
\newtheorem{exercise}{Exercise}
\newtheorem*{remark}{Remark}

\newenvironment{solution}{\par\noindent\textbf{\color{solFR}Solution.}\,}{\hfill\color{solFR}$\blacksquare$\par}

% tcolorbox styling for theorem-like environments.
\tcolorboxenvironment{definition}{enhanced,breakable,colback=defBG,colframe=defFR,boxrule=0.6pt,arc=2pt,left=3mm,right=3mm,top=1.5mm,bottom=1.5mm,before skip=5pt,after skip=5pt}
\tcolorboxenvironment{theorem}{enhanced,breakable,colback=thmBG,colframe=thmFR,boxrule=0.6pt,arc=2pt,left=3mm,right=3mm,top=1.5mm,bottom=1.5mm,before skip=5pt,after skip=5pt}
\tcolorboxenvironment{lemma}{enhanced,breakable,colback=thmBG,colframe=thmFR,boxrule=0.6pt,arc=2pt,left=3mm,right=3mm,top=1.5mm,bottom=1.5mm,before skip=5pt,after skip=5pt}
\tcolorboxenvironment{example}{enhanced,breakable,colback=exaBG,colframe=exaFR,leftrule=3pt,rightrule=0pt,toprule=0pt,bottomrule=0pt,arc=0pt,left=3mm,right=2mm,top=1mm,bottom=1mm,before skip=4pt,after skip=4pt}
\tcolorboxenvironment{proof}{enhanced,breakable,colback=proofBG,colframe=proofFR,leftrule=2pt,rightrule=0pt,toprule=0pt,bottomrule=0pt,arc=0pt,left=3mm,right=2mm,top=1mm,bottom=1mm,before skip=4pt,after skip=4pt}
\tcolorboxenvironment{exercise}{enhanced,breakable,colback=exBG,colframe=exFR,leftrule=3pt,rightrule=0pt,toprule=0pt,bottomrule=0pt,arc=0pt,left=3mm,right=2mm,top=1mm,bottom=1mm,before skip=4pt,after skip=4pt}
\tcolorboxenvironment{solution}{enhanced,breakable,colback=solBG,colframe=solFR,leftrule=3pt,rightrule=0pt,toprule=0pt,bottomrule=0pt,arc=0pt,left=3mm,right=2mm,top=1mm,bottom=1mm,before skip=3pt,after skip=6pt}
\tcolorboxenvironment{remark}{enhanced,colback=rmkBG,colframe=rmkFR,leftrule=2pt,rightrule=0pt,toprule=0pt,bottomrule=0pt,arc=0pt,left=3mm,right=2mm,top=1mm,bottom=1mm,before skip=4pt,after skip=4pt}

% Centered tabular block. One optional argument:
%   #1 = scale factor for \scalebox (positive decimal, default 1.0), OR
%        a font-size command (\small, \footnotesize, ...) to shrink the body's
%        font.
% Vertical space before/after is fixed at 2mm.
% Usage:
%   \begin{ctab} ... \end{ctab}              % no scaling, no font change
%   \begin{ctab}[0.85] ... \end{ctab}        % \scalebox at 0.85
%   \begin{ctab}[\small] ... \end{ctab}      % \small font inside the table
% Replaces the manual `\vspace{2mm}{\centering ... \par}\vspace{2mm}` pattern
% (CONVENTIONS.md §9). The lead-in prose paragraph must end before \begin{ctab}
% (a blank line, or a paragraph-ending statement) so the lead-in is not centered.
% Note: with a numeric scale, \scalebox wraps the body in a single hbox at any
% scale, so a ctab-wrapped tabular cannot break across pages. For long tables
% that must break, fall back to the manual `{\centering ... \par}` pattern
% (CONVENTIONS.md §9) instead of ctab.
\newsavebox{\ctabsavebox}
\newenvironment{ctab}[1][1.0]
  {\def\ctabarg{#1}\edef\ctabargstr{\detokenize{#1}}%
   \par\vspace{2mm}\begin{lrbox}{\ctabsavebox}%
   \IfDecimal{\ctabargstr}{}{#1}\ignorespaces}
  {\end{lrbox}\noindent{\centering
    \IfDecimal{\ctabargstr}%
      {\scalebox{\ctabarg}{\usebox{\ctabsavebox}}}%
      {\usebox{\ctabsavebox}}\par}%
   \vspace{2mm}}

% Emphasis: bold rather than italic (italic is hard to read on screen).
\renewcommand{\emph}[1]{\textbf{#1}}

% Proof heading: bold (override amsthm's default \itshape).
\makeatletter
\renewenvironment{proof}[1][\proofname]{\par
  \pushQED{\qed}%
  \normalfont \topsep6\p@\@plus6\p@\relax
  \trivlist
  \item[\hskip\labelsep
        \bfseries
    #1\@addpunct{.}]\ignorespaces
}{%
  \popQED\endtrivlist\@endpefalse
}
\makeatother

% Three signature pedagogical boxes:
%   motivation (blue)  : opening hooks, contextual framing
%   derivation (green) : supplementary derivations beyond the main flow
%   suppremark (orange): instructor commentary, study tips, broader perspective
\newtcolorbox{motivation}[1][Motivation]{breakable, enhanced, colback=blue!5,
  colframe=blue!50!black, coltitle=white,
  title=\textbf{#1}, fonttitle=\bfseries,
  attach boxed title to top left={yshift=-2mm,xshift=3mm},
  boxed title style={colback=blue!50!black}}
\newtcolorbox{derivation}[1][Derivation]{breakable, enhanced, colback=green!5,
  colframe=green!40!black, coltitle=white,
  title=\textbf{#1}, fonttitle=\bfseries,
  attach boxed title to top left={yshift=-2mm,xshift=3mm},
  boxed title style={colback=green!40!black}}
\newtcolorbox{suppremark}[1][Remark]{breakable, enhanced, colback=orange!5,
  colframe=orange!60!black, coltitle=white,
  title=\textbf{#1}, fonttitle=\bfseries,
  attach boxed title to top left={yshift=-2mm,xshift=3mm},
  boxed title style={colback=orange!60!black}}

% Custom math operators used throughout the course.
\usepackage{dsfont}
\DeclareMathOperator\prb{\mathsf{P}}
\DeclareMathOperator\expc{\mathsf{E}}
\DeclareMathOperator\var{var}
\DeclareMathOperator\cov{cov}
\DeclareMathOperator\cor{cor}
\DeclareMathOperator\indc{\mathds{1}}
\DeclareMathOperator\dom{dom}
\DeclareMathOperator\epi{epi}
\DeclareMathOperator\diag{diag}
\DeclareMathOperator\rk{rank}
\DeclareMathOperator\nul{null}
\DeclareMathOperator\col{col}
\DeclareMathOperator\img{img}
\let\ker\relax
\DeclareMathOperator\ker{ker}
\DeclareMathOperator\tr{tr}
\DeclareMathOperator\df{df}

\newcommand{\dd}{\mathrm{d}}

\usepackage{parskip}
\setlength{\parskip}{3pt}\setlength{\jot}{4pt}

\usepackage{titlesec}
\titleformat{\section}{\color{hdrBG}\normalfont\Large\bfseries}{\thesection}{0.7em}{}
\titleformat{\subsection}{\color{accent}\normalfont\large\bfseries}{\thesubsection}{0.6em}{}
\titlespacing*{\section}{0pt}{4pt}{2pt}
\titlespacing*{\subsection}{0pt}{4pt}{1pt}

%  immediately after \documentclass[17pt]{extarticle}.
%
%  Unit-specific configuration that must remain in lecture.Rnw, after this
%  \input but before \begin{document}:
%    \fancyhead[L]{\small\color{hdrBG}\textbf{Unit\ N <chapter title>}}
%    \setcounter{theorem}{0}\setcounter{example}{0}
%    \renewcommand\thedefinition{N.\arabic{theorem}}
%    \renewcommand\thetheorem  {N.\arabic{theorem}}
%    \renewcommand\theexample  {N.\arabic{example}}
%    \renewcommand\theexercise {N.\arabic{exercise}}
%    \renewcommand\thefigure   {N.\arabic{figure}}
%    \title{\color{hdrBG}\bfseries\Huge <chapter title>}\author{}\date{}
% =============================================================================

\usepackage[
  paperwidth=257mm,
  paperheight=144.5mm,
  top=9mm,bottom=9mm,left=6mm,right=6mm,
  headheight=7mm,headsep=2mm,footskip=19pt
]{geometry}

\usepackage{amsmath,amsthm,amssymb,amsfonts}
\usepackage{xstring}  % \IfDecimal used by ctab for numeric-vs-font dispatch
\usepackage{booktabs,array,longtable,tabularx}
\usepackage[shortlabels]{enumitem}
\setlist[itemize,1]{leftmargin=1.5em, labelsep=0.35em}
\usepackage{mathtools}
\usepackage{xcolor}
\usepackage{fancyhdr}
\usepackage{tikz}
\usetikzlibrary{arrows.meta,calc,decorations.markings}
\usepackage[most]{tcolorbox}
\tcbuselibrary{theorems,skins,breakable}
\usepackage[round,authoryear]{natbib}  % loaded before hyperref so \citep is hyperlinked
\usepackage{hyperref}
\hypersetup{colorlinks=true,linkcolor=blue!60!black,urlcolor=blue!70!black,citecolor=blue!60!black}

% Source code highlighting (requires pygmentize; compile with -shell-escape)
\usepackage[cache=false]{minted}
\setminted{fontsize=\small, frame=single, framesep=2mm, rulecolor=gray!40,
  bgcolor=gray!5, breaklines=false, tabsize=2,
  autogobble=true,
  linenos=true,
  numbersep=5pt,
  xleftmargin=2.5em}

% Line-number style: gray, slightly smaller, so it does not compete with the code.
\renewcommand{\theFancyVerbLine}{%
  \textcolor{gray!60}{\scriptsize\arabic{FancyVerbLine}}}

% xeCJK is loaded for safety even though all course content is in English;
% it lets occasional CJK characters (e.g. a student name in a sample dataset)
% render without breaking the build.
\usepackage{xeCJK}
\setCJKmainfont{Noto Serif CJK TC}
\setCJKsansfont{Noto Sans CJK TC}
\setCJKmonofont{Noto Sans Mono CJK TC}

% Course-wide color palette.
\definecolor{hdrBG}{HTML}{1E3A8A}
\definecolor{accent}{HTML}{0EA5E9}
\definecolor{defBG}{HTML}{EFF6FF}\definecolor{defFR}{HTML}{2563EB}
\definecolor{thmBG}{HTML}{FEF3C7}\definecolor{thmFR}{HTML}{B45309}
\definecolor{exaBG}{HTML}{FEF2F2}\definecolor{exaFR}{HTML}{B91C1C}
\definecolor{proofBG}{HTML}{F8FAFC}\definecolor{proofFR}{HTML}{64748B}
\definecolor{exBG}{HTML}{FFF7ED}\definecolor{exFR}{HTML}{C2410C}
\definecolor{solBG}{HTML}{F0FDF4}\definecolor{solFR}{HTML}{15803D}
\definecolor{rmkBG}{HTML}{F1F5F9}\definecolor{rmkFR}{HTML}{94A3B8}

% Header / footer skeleton. \fancyhead[L] is set per-unit in lecture.Rnw.
\pagestyle{fancy}\fancyhf{}
\renewcommand{\headrulewidth}{0.4pt}
\renewcommand{\headrule}{\color{hdrBG!30}\hrule width\textwidth height0.4pt}
\fancyhead[R]{\small\color{hdrBG}\leftmark}
\fancyfoot[R]{\small\color{gray}\thepage}
\renewcommand{\sectionmark}[1]{\markboth{\S\thesection~#1}{}}

% Theorem-like environments. All use \theoremstyle{definition} so the body
% renders upright (no italic / slanted text in theorem statements, proofs, or
% solutions).
\theoremstyle{definition}
\newtheorem{theorem}{Theorem}
\newtheorem{lemma}[theorem]{Lemma}
\newtheorem{definition}[theorem]{Definition}
\newtheorem{example}{Example}
\newtheorem{exercise}{Exercise}
\newtheorem*{remark}{Remark}

\newenvironment{solution}{\par\noindent\textbf{\color{solFR}Solution.}\,}{\hfill\color{solFR}$\blacksquare$\par}

% tcolorbox styling for theorem-like environments.
\tcolorboxenvironment{definition}{enhanced,breakable,colback=defBG,colframe=defFR,boxrule=0.6pt,arc=2pt,left=3mm,right=3mm,top=1.5mm,bottom=1.5mm,before skip=5pt,after skip=5pt}
\tcolorboxenvironment{theorem}{enhanced,breakable,colback=thmBG,colframe=thmFR,boxrule=0.6pt,arc=2pt,left=3mm,right=3mm,top=1.5mm,bottom=1.5mm,before skip=5pt,after skip=5pt}
\tcolorboxenvironment{lemma}{enhanced,breakable,colback=thmBG,colframe=thmFR,boxrule=0.6pt,arc=2pt,left=3mm,right=3mm,top=1.5mm,bottom=1.5mm,before skip=5pt,after skip=5pt}
\tcolorboxenvironment{example}{enhanced,breakable,colback=exaBG,colframe=exaFR,leftrule=3pt,rightrule=0pt,toprule=0pt,bottomrule=0pt,arc=0pt,left=3mm,right=2mm,top=1mm,bottom=1mm,before skip=4pt,after skip=4pt}
\tcolorboxenvironment{proof}{enhanced,breakable,colback=proofBG,colframe=proofFR,leftrule=2pt,rightrule=0pt,toprule=0pt,bottomrule=0pt,arc=0pt,left=3mm,right=2mm,top=1mm,bottom=1mm,before skip=4pt,after skip=4pt}
\tcolorboxenvironment{exercise}{enhanced,breakable,colback=exBG,colframe=exFR,leftrule=3pt,rightrule=0pt,toprule=0pt,bottomrule=0pt,arc=0pt,left=3mm,right=2mm,top=1mm,bottom=1mm,before skip=4pt,after skip=4pt}
\tcolorboxenvironment{solution}{enhanced,breakable,colback=solBG,colframe=solFR,leftrule=3pt,rightrule=0pt,toprule=0pt,bottomrule=0pt,arc=0pt,left=3mm,right=2mm,top=1mm,bottom=1mm,before skip=3pt,after skip=6pt}
\tcolorboxenvironment{remark}{enhanced,colback=rmkBG,colframe=rmkFR,leftrule=2pt,rightrule=0pt,toprule=0pt,bottomrule=0pt,arc=0pt,left=3mm,right=2mm,top=1mm,bottom=1mm,before skip=4pt,after skip=4pt}

% Centered tabular block. One optional argument:
%   #1 = scale factor for \scalebox (positive decimal, default 1.0), OR
%        a font-size command (\small, \footnotesize, ...) to shrink the body's
%        font.
% Vertical space before/after is fixed at 2mm.
% Usage:
%   \begin{ctab} ... \end{ctab}              % no scaling, no font change
%   \begin{ctab}[0.85] ... \end{ctab}        % \scalebox at 0.85
%   \begin{ctab}[\small] ... \end{ctab}      % \small font inside the table
% Replaces the manual `\vspace{2mm}{\centering ... \par}\vspace{2mm}` pattern
% (CONVENTIONS.md §9). The lead-in prose paragraph must end before \begin{ctab}
% (a blank line, or a paragraph-ending statement) so the lead-in is not centered.
% Note: with a numeric scale, \scalebox wraps the body in a single hbox at any
% scale, so a ctab-wrapped tabular cannot break across pages. For long tables
% that must break, fall back to the manual `{\centering ... \par}` pattern
% (CONVENTIONS.md §9) instead of ctab.
\newsavebox{\ctabsavebox}
\newenvironment{ctab}[1][1.0]
  {\def\ctabarg{#1}\edef\ctabargstr{\detokenize{#1}}%
   \par\vspace{2mm}\begin{lrbox}{\ctabsavebox}%
   \IfDecimal{\ctabargstr}{}{#1}\ignorespaces}
  {\end{lrbox}\noindent{\centering
    \IfDecimal{\ctabargstr}%
      {\scalebox{\ctabarg}{\usebox{\ctabsavebox}}}%
      {\usebox{\ctabsavebox}}\par}%
   \vspace{2mm}}

% Emphasis: bold rather than italic (italic is hard to read on screen).
\renewcommand{\emph}[1]{\textbf{#1}}

% Proof heading: bold (override amsthm's default \itshape).
\makeatletter
\renewenvironment{proof}[1][\proofname]{\par
  \pushQED{\qed}%
  \normalfont \topsep6\p@\@plus6\p@\relax
  \trivlist
  \item[\hskip\labelsep
        \bfseries
    #1\@addpunct{.}]\ignorespaces
}{%
  \popQED\endtrivlist\@endpefalse
}
\makeatother

% Three signature pedagogical boxes:
%   motivation (blue)  : opening hooks, contextual framing
%   derivation (green) : supplementary derivations beyond the main flow
%   suppremark (orange): instructor commentary, study tips, broader perspective
\newtcolorbox{motivation}[1][Motivation]{breakable, enhanced, colback=blue!5,
  colframe=blue!50!black, coltitle=white,
  title=\textbf{#1}, fonttitle=\bfseries,
  attach boxed title to top left={yshift=-2mm,xshift=3mm},
  boxed title style={colback=blue!50!black}}
\newtcolorbox{derivation}[1][Derivation]{breakable, enhanced, colback=green!5,
  colframe=green!40!black, coltitle=white,
  title=\textbf{#1}, fonttitle=\bfseries,
  attach boxed title to top left={yshift=-2mm,xshift=3mm},
  boxed title style={colback=green!40!black}}
\newtcolorbox{suppremark}[1][Remark]{breakable, enhanced, colback=orange!5,
  colframe=orange!60!black, coltitle=white,
  title=\textbf{#1}, fonttitle=\bfseries,
  attach boxed title to top left={yshift=-2mm,xshift=3mm},
  boxed title style={colback=orange!60!black}}

% Custom math operators used throughout the course.
\usepackage{dsfont}
\DeclareMathOperator\prb{\mathsf{P}}
\DeclareMathOperator\expc{\mathsf{E}}
\DeclareMathOperator\var{var}
\DeclareMathOperator\cov{cov}
\DeclareMathOperator\cor{cor}
\DeclareMathOperator\indc{\mathds{1}}
\DeclareMathOperator\dom{dom}
\DeclareMathOperator\epi{epi}
\DeclareMathOperator\diag{diag}
\DeclareMathOperator\rk{rank}
\DeclareMathOperator\nul{null}
\DeclareMathOperator\col{col}
\DeclareMathOperator\img{img}
\let\ker\relax
\DeclareMathOperator\ker{ker}
\DeclareMathOperator\tr{tr}
\DeclareMathOperator\df{df}

\newcommand{\dd}{\mathrm{d}}

\usepackage{parskip}
\setlength{\parskip}{3pt}\setlength{\jot}{4pt}

\usepackage{titlesec}
\titleformat{\section}{\color{hdrBG}\normalfont\Large\bfseries}{\thesection}{0.7em}{}
\titleformat{\subsection}{\color{accent}\normalfont\large\bfseries}{\thesubsection}{0.6em}{}
\titlespacing*{\section}{0pt}{4pt}{2pt}
\titlespacing*{\subsection}{0pt}{4pt}{1pt}

%  immediately after \documentclass[17pt]{extarticle}.
%
%  Unit-specific configuration that must remain in lecture.Rnw, after this
%  \input but before \begin{document}:
%    \fancyhead[L]{\small\color{hdrBG}\textbf{Unit\ N <chapter title>}}
%    \setcounter{theorem}{0}\setcounter{example}{0}
%    \renewcommand\thedefinition{N.\arabic{theorem}}
%    \renewcommand\thetheorem  {N.\arabic{theorem}}
%    \renewcommand\theexample  {N.\arabic{example}}
%    \renewcommand\theexercise {N.\arabic{exercise}}
%    \renewcommand\thefigure   {N.\arabic{figure}}
%    \title{\color{hdrBG}\bfseries\Huge <chapter title>}\author{}\date{}
% =============================================================================

\usepackage[
  paperwidth=257mm,
  paperheight=144.5mm,
  top=9mm,bottom=9mm,left=6mm,right=6mm,
  headheight=7mm,headsep=2mm,footskip=19pt
]{geometry}

\usepackage{amsmath,amsthm,amssymb,amsfonts}
\usepackage{xstring}  % \IfDecimal used by ctab for numeric-vs-font dispatch
\usepackage{booktabs,array,longtable,tabularx}
\usepackage[shortlabels]{enumitem}
\setlist[itemize,1]{leftmargin=1.5em, labelsep=0.35em}
\usepackage{mathtools}
\usepackage{xcolor}
\usepackage{fancyhdr}
\usepackage{tikz}
\usetikzlibrary{arrows.meta,calc,decorations.markings}
\usepackage[most]{tcolorbox}
\tcbuselibrary{theorems,skins,breakable}
\usepackage[round,authoryear]{natbib}  % loaded before hyperref so \citep is hyperlinked
\usepackage{hyperref}
\hypersetup{colorlinks=true,linkcolor=blue!60!black,urlcolor=blue!70!black,citecolor=blue!60!black}

% Source code highlighting (requires pygmentize; compile with -shell-escape)
\usepackage[cache=false]{minted}
\setminted{fontsize=\small, frame=single, framesep=2mm, rulecolor=gray!40,
  bgcolor=gray!5, breaklines=false, tabsize=2,
  autogobble=true,
  linenos=true,
  numbersep=5pt,
  xleftmargin=2.5em}

% Line-number style: gray, slightly smaller, so it does not compete with the code.
\renewcommand{\theFancyVerbLine}{%
  \textcolor{gray!60}{\scriptsize\arabic{FancyVerbLine}}}

% xeCJK is loaded for safety even though all course content is in English;
% it lets occasional CJK characters (e.g. a student name in a sample dataset)
% render without breaking the build.
\usepackage{xeCJK}
\setCJKmainfont{Noto Serif CJK TC}
\setCJKsansfont{Noto Sans CJK TC}
\setCJKmonofont{Noto Sans Mono CJK TC}

% Course-wide color palette.
\definecolor{hdrBG}{HTML}{1E3A8A}
\definecolor{accent}{HTML}{0EA5E9}
\definecolor{defBG}{HTML}{EFF6FF}\definecolor{defFR}{HTML}{2563EB}
\definecolor{thmBG}{HTML}{FEF3C7}\definecolor{thmFR}{HTML}{B45309}
\definecolor{exaBG}{HTML}{FEF2F2}\definecolor{exaFR}{HTML}{B91C1C}
\definecolor{proofBG}{HTML}{F8FAFC}\definecolor{proofFR}{HTML}{64748B}
\definecolor{exBG}{HTML}{FFF7ED}\definecolor{exFR}{HTML}{C2410C}
\definecolor{solBG}{HTML}{F0FDF4}\definecolor{solFR}{HTML}{15803D}
\definecolor{rmkBG}{HTML}{F1F5F9}\definecolor{rmkFR}{HTML}{94A3B8}

% Header / footer skeleton. \fancyhead[L] is set per-unit in lecture.Rnw.
\pagestyle{fancy}\fancyhf{}
\renewcommand{\headrulewidth}{0.4pt}
\renewcommand{\headrule}{\color{hdrBG!30}\hrule width\textwidth height0.4pt}
\fancyhead[R]{\small\color{hdrBG}\leftmark}
\fancyfoot[R]{\small\color{gray}\thepage}
\renewcommand{\sectionmark}[1]{\markboth{\S\thesection~#1}{}}

% Theorem-like environments. All use \theoremstyle{definition} so the body
% renders upright (no italic / slanted text in theorem statements, proofs, or
% solutions).
\theoremstyle{definition}
\newtheorem{theorem}{Theorem}
\newtheorem{lemma}[theorem]{Lemma}
\newtheorem{definition}[theorem]{Definition}
\newtheorem{example}{Example}
\newtheorem{exercise}{Exercise}
\newtheorem*{remark}{Remark}

\newenvironment{solution}{\par\noindent\textbf{\color{solFR}Solution.}\,}{\hfill\color{solFR}$\blacksquare$\par}

% tcolorbox styling for theorem-like environments.
\tcolorboxenvironment{definition}{enhanced,breakable,colback=defBG,colframe=defFR,boxrule=0.6pt,arc=2pt,left=3mm,right=3mm,top=1.5mm,bottom=1.5mm,before skip=5pt,after skip=5pt}
\tcolorboxenvironment{theorem}{enhanced,breakable,colback=thmBG,colframe=thmFR,boxrule=0.6pt,arc=2pt,left=3mm,right=3mm,top=1.5mm,bottom=1.5mm,before skip=5pt,after skip=5pt}
\tcolorboxenvironment{lemma}{enhanced,breakable,colback=thmBG,colframe=thmFR,boxrule=0.6pt,arc=2pt,left=3mm,right=3mm,top=1.5mm,bottom=1.5mm,before skip=5pt,after skip=5pt}
\tcolorboxenvironment{example}{enhanced,breakable,colback=exaBG,colframe=exaFR,leftrule=3pt,rightrule=0pt,toprule=0pt,bottomrule=0pt,arc=0pt,left=3mm,right=2mm,top=1mm,bottom=1mm,before skip=4pt,after skip=4pt}
\tcolorboxenvironment{proof}{enhanced,breakable,colback=proofBG,colframe=proofFR,leftrule=2pt,rightrule=0pt,toprule=0pt,bottomrule=0pt,arc=0pt,left=3mm,right=2mm,top=1mm,bottom=1mm,before skip=4pt,after skip=4pt}
\tcolorboxenvironment{exercise}{enhanced,breakable,colback=exBG,colframe=exFR,leftrule=3pt,rightrule=0pt,toprule=0pt,bottomrule=0pt,arc=0pt,left=3mm,right=2mm,top=1mm,bottom=1mm,before skip=4pt,after skip=4pt}
\tcolorboxenvironment{solution}{enhanced,breakable,colback=solBG,colframe=solFR,leftrule=3pt,rightrule=0pt,toprule=0pt,bottomrule=0pt,arc=0pt,left=3mm,right=2mm,top=1mm,bottom=1mm,before skip=3pt,after skip=6pt}
\tcolorboxenvironment{remark}{enhanced,colback=rmkBG,colframe=rmkFR,leftrule=2pt,rightrule=0pt,toprule=0pt,bottomrule=0pt,arc=0pt,left=3mm,right=2mm,top=1mm,bottom=1mm,before skip=4pt,after skip=4pt}

% Centered tabular block. One optional argument:
%   #1 = scale factor for \scalebox (positive decimal, default 1.0), OR
%        a font-size command (\small, \footnotesize, ...) to shrink the body's
%        font.
% Vertical space before/after is fixed at 2mm.
% Usage:
%   \begin{ctab} ... \end{ctab}              % no scaling, no font change
%   \begin{ctab}[0.85] ... \end{ctab}        % \scalebox at 0.85
%   \begin{ctab}[\small] ... \end{ctab}      % \small font inside the table
% Replaces the manual `\vspace{2mm}{\centering ... \par}\vspace{2mm}` pattern
% (CONVENTIONS.md §9). The lead-in prose paragraph must end before \begin{ctab}
% (a blank line, or a paragraph-ending statement) so the lead-in is not centered.
% Note: with a numeric scale, \scalebox wraps the body in a single hbox at any
% scale, so a ctab-wrapped tabular cannot break across pages. For long tables
% that must break, fall back to the manual `{\centering ... \par}` pattern
% (CONVENTIONS.md §9) instead of ctab.
\newsavebox{\ctabsavebox}
\newenvironment{ctab}[1][1.0]
  {\def\ctabarg{#1}\edef\ctabargstr{\detokenize{#1}}%
   \par\vspace{2mm}\begin{lrbox}{\ctabsavebox}%
   \IfDecimal{\ctabargstr}{}{#1}\ignorespaces}
  {\end{lrbox}\noindent{\centering
    \IfDecimal{\ctabargstr}%
      {\scalebox{\ctabarg}{\usebox{\ctabsavebox}}}%
      {\usebox{\ctabsavebox}}\par}%
   \vspace{2mm}}

% Emphasis: bold rather than italic (italic is hard to read on screen).
\renewcommand{\emph}[1]{\textbf{#1}}

% Proof heading: bold (override amsthm's default \itshape).
\makeatletter
\renewenvironment{proof}[1][\proofname]{\par
  \pushQED{\qed}%
  \normalfont \topsep6\p@\@plus6\p@\relax
  \trivlist
  \item[\hskip\labelsep
        \bfseries
    #1\@addpunct{.}]\ignorespaces
}{%
  \popQED\endtrivlist\@endpefalse
}
\makeatother

% Three signature pedagogical boxes:
%   motivation (blue)  : opening hooks, contextual framing
%   derivation (green) : supplementary derivations beyond the main flow
%   suppremark (orange): instructor commentary, study tips, broader perspective
\newtcolorbox{motivation}[1][Motivation]{breakable, enhanced, colback=blue!5,
  colframe=blue!50!black, coltitle=white,
  title=\textbf{#1}, fonttitle=\bfseries,
  attach boxed title to top left={yshift=-2mm,xshift=3mm},
  boxed title style={colback=blue!50!black}}
\newtcolorbox{derivation}[1][Derivation]{breakable, enhanced, colback=green!5,
  colframe=green!40!black, coltitle=white,
  title=\textbf{#1}, fonttitle=\bfseries,
  attach boxed title to top left={yshift=-2mm,xshift=3mm},
  boxed title style={colback=green!40!black}}
\newtcolorbox{suppremark}[1][Remark]{breakable, enhanced, colback=orange!5,
  colframe=orange!60!black, coltitle=white,
  title=\textbf{#1}, fonttitle=\bfseries,
  attach boxed title to top left={yshift=-2mm,xshift=3mm},
  boxed title style={colback=orange!60!black}}

% Custom math operators used throughout the course.
\usepackage{dsfont}
\DeclareMathOperator\prb{\mathsf{P}}
\DeclareMathOperator\expc{\mathsf{E}}
\DeclareMathOperator\var{var}
\DeclareMathOperator\cov{cov}
\DeclareMathOperator\cor{cor}
\DeclareMathOperator\indc{\mathds{1}}
\DeclareMathOperator\dom{dom}
\DeclareMathOperator\epi{epi}
\DeclareMathOperator\diag{diag}
\DeclareMathOperator\rk{rank}
\DeclareMathOperator\nul{null}
\DeclareMathOperator\col{col}
\DeclareMathOperator\img{img}
\let\ker\relax
\DeclareMathOperator\ker{ker}
\DeclareMathOperator\tr{tr}
\DeclareMathOperator\df{df}

\newcommand{\dd}{\mathrm{d}}

\usepackage{parskip}
\setlength{\parskip}{3pt}\setlength{\jot}{4pt}

\usepackage{titlesec}
\titleformat{\section}{\color{hdrBG}\normalfont\Large\bfseries}{\thesection}{0.7em}{}
\titleformat{\subsection}{\color{accent}\normalfont\large\bfseries}{\thesubsection}{0.6em}{}
\titlespacing*{\section}{0pt}{4pt}{2pt}
\titlespacing*{\subsection}{0pt}{4pt}{1pt}

%  immediately after \documentclass[17pt]{extarticle}.
%
%  Unit-specific configuration that must remain in lecture.Rnw, after this
%  \input but before \begin{document}:
%    \fancyhead[L]{\small\color{hdrBG}\textbf{Unit\ N <chapter title>}}
%    \setcounter{theorem}{0}\setcounter{example}{0}
%    \renewcommand\thedefinition{N.\arabic{theorem}}
%    \renewcommand\thetheorem  {N.\arabic{theorem}}
%    \renewcommand\theexample  {N.\arabic{example}}
%    \renewcommand\theexercise {N.\arabic{exercise}}
%    \renewcommand\thefigure   {N.\arabic{figure}}
%    \title{\color{hdrBG}\bfseries\Huge <chapter title>}\author{}\date{}
% =============================================================================

\usepackage[
  paperwidth=257mm,
  paperheight=144.5mm,
  top=9mm,bottom=9mm,left=6mm,right=6mm,
  headheight=7mm,headsep=2mm,footskip=19pt
]{geometry}

\usepackage{amsmath,amsthm,amssymb,amsfonts}
\usepackage{xstring}  % \IfDecimal used by ctab for numeric-vs-font dispatch
\usepackage{booktabs,array,longtable,tabularx}
\usepackage[shortlabels]{enumitem}
\setlist[itemize,1]{leftmargin=1.5em, labelsep=0.35em}
\usepackage{mathtools}
\usepackage{xcolor}
\usepackage{fancyhdr}
\usepackage{tikz}
\usetikzlibrary{arrows.meta,calc,decorations.markings}
\usepackage[most]{tcolorbox}
\tcbuselibrary{theorems,skins,breakable}
\usepackage[round,authoryear]{natbib}  % loaded before hyperref so \citep is hyperlinked
\usepackage{hyperref}
\hypersetup{colorlinks=true,linkcolor=blue!60!black,urlcolor=blue!70!black,citecolor=blue!60!black}

% Source code highlighting (requires pygmentize; compile with -shell-escape)
\usepackage[cache=false]{minted}
\setminted{fontsize=\small, frame=single, framesep=2mm, rulecolor=gray!40,
  bgcolor=gray!5, breaklines=false, tabsize=2,
  autogobble=true,
  linenos=true,
  numbersep=5pt,
  xleftmargin=2.5em}

% Line-number style: gray, slightly smaller, so it does not compete with the code.
\renewcommand{\theFancyVerbLine}{%
  \textcolor{gray!60}{\scriptsize\arabic{FancyVerbLine}}}

% xeCJK is loaded for safety even though all course content is in English;
% it lets occasional CJK characters (e.g. a student name in a sample dataset)
% render without breaking the build.
\usepackage{xeCJK}
\setCJKmainfont{Noto Serif CJK TC}
\setCJKsansfont{Noto Sans CJK TC}
\setCJKmonofont{Noto Sans Mono CJK TC}

% Course-wide color palette.
\definecolor{hdrBG}{HTML}{1E3A8A}
\definecolor{accent}{HTML}{0EA5E9}
\definecolor{defBG}{HTML}{EFF6FF}\definecolor{defFR}{HTML}{2563EB}
\definecolor{thmBG}{HTML}{FEF3C7}\definecolor{thmFR}{HTML}{B45309}
\definecolor{exaBG}{HTML}{FEF2F2}\definecolor{exaFR}{HTML}{B91C1C}
\definecolor{proofBG}{HTML}{F8FAFC}\definecolor{proofFR}{HTML}{64748B}
\definecolor{exBG}{HTML}{FFF7ED}\definecolor{exFR}{HTML}{C2410C}
\definecolor{solBG}{HTML}{F0FDF4}\definecolor{solFR}{HTML}{15803D}
\definecolor{rmkBG}{HTML}{F1F5F9}\definecolor{rmkFR}{HTML}{94A3B8}

% Header / footer skeleton. \fancyhead[L] is set per-unit in lecture.Rnw.
\pagestyle{fancy}\fancyhf{}
\renewcommand{\headrulewidth}{0.4pt}
\renewcommand{\headrule}{\color{hdrBG!30}\hrule width\textwidth height0.4pt}
\fancyhead[R]{\small\color{hdrBG}\leftmark}
\fancyfoot[R]{\small\color{gray}\thepage}
\renewcommand{\sectionmark}[1]{\markboth{\S\thesection~#1}{}}

% Theorem-like environments. All use \theoremstyle{definition} so the body
% renders upright (no italic / slanted text in theorem statements, proofs, or
% solutions).
\theoremstyle{definition}
\newtheorem{theorem}{Theorem}
\newtheorem{lemma}[theorem]{Lemma}
\newtheorem{definition}[theorem]{Definition}
\newtheorem{example}{Example}
\newtheorem{exercise}{Exercise}
\newtheorem*{remark}{Remark}

\newenvironment{solution}{\par\noindent\textbf{\color{solFR}Solution.}\,}{\hfill\color{solFR}$\blacksquare$\par}

% tcolorbox styling for theorem-like environments.
\tcolorboxenvironment{definition}{enhanced,breakable,colback=defBG,colframe=defFR,boxrule=0.6pt,arc=2pt,left=3mm,right=3mm,top=1.5mm,bottom=1.5mm,before skip=5pt,after skip=5pt}
\tcolorboxenvironment{theorem}{enhanced,breakable,colback=thmBG,colframe=thmFR,boxrule=0.6pt,arc=2pt,left=3mm,right=3mm,top=1.5mm,bottom=1.5mm,before skip=5pt,after skip=5pt}
\tcolorboxenvironment{lemma}{enhanced,breakable,colback=thmBG,colframe=thmFR,boxrule=0.6pt,arc=2pt,left=3mm,right=3mm,top=1.5mm,bottom=1.5mm,before skip=5pt,after skip=5pt}
\tcolorboxenvironment{example}{enhanced,breakable,colback=exaBG,colframe=exaFR,leftrule=3pt,rightrule=0pt,toprule=0pt,bottomrule=0pt,arc=0pt,left=3mm,right=2mm,top=1mm,bottom=1mm,before skip=4pt,after skip=4pt}
\tcolorboxenvironment{proof}{enhanced,breakable,colback=proofBG,colframe=proofFR,leftrule=2pt,rightrule=0pt,toprule=0pt,bottomrule=0pt,arc=0pt,left=3mm,right=2mm,top=1mm,bottom=1mm,before skip=4pt,after skip=4pt}
\tcolorboxenvironment{exercise}{enhanced,breakable,colback=exBG,colframe=exFR,leftrule=3pt,rightrule=0pt,toprule=0pt,bottomrule=0pt,arc=0pt,left=3mm,right=2mm,top=1mm,bottom=1mm,before skip=4pt,after skip=4pt}
\tcolorboxenvironment{solution}{enhanced,breakable,colback=solBG,colframe=solFR,leftrule=3pt,rightrule=0pt,toprule=0pt,bottomrule=0pt,arc=0pt,left=3mm,right=2mm,top=1mm,bottom=1mm,before skip=3pt,after skip=6pt}
\tcolorboxenvironment{remark}{enhanced,colback=rmkBG,colframe=rmkFR,leftrule=2pt,rightrule=0pt,toprule=0pt,bottomrule=0pt,arc=0pt,left=3mm,right=2mm,top=1mm,bottom=1mm,before skip=4pt,after skip=4pt}

% Centered tabular block. One optional argument:
%   #1 = scale factor for \scalebox (positive decimal, default 1.0), OR
%        a font-size command (\small, \footnotesize, ...) to shrink the body's
%        font.
% Vertical space before/after is fixed at 2mm.
% Usage:
%   \begin{ctab} ... \end{ctab}              % no scaling, no font change
%   \begin{ctab}[0.85] ... \end{ctab}        % \scalebox at 0.85
%   \begin{ctab}[\small] ... \end{ctab}      % \small font inside the table
% Replaces the manual `\vspace{2mm}{\centering ... \par}\vspace{2mm}` pattern
% (CONVENTIONS.md §9). The lead-in prose paragraph must end before \begin{ctab}
% (a blank line, or a paragraph-ending statement) so the lead-in is not centered.
% Note: with a numeric scale, \scalebox wraps the body in a single hbox at any
% scale, so a ctab-wrapped tabular cannot break across pages. For long tables
% that must break, fall back to the manual `{\centering ... \par}` pattern
% (CONVENTIONS.md §9) instead of ctab.
\newsavebox{\ctabsavebox}
\newenvironment{ctab}[1][1.0]
  {\def\ctabarg{#1}\edef\ctabargstr{\detokenize{#1}}%
   \par\vspace{2mm}\begin{lrbox}{\ctabsavebox}%
   \IfDecimal{\ctabargstr}{}{#1}\ignorespaces}
  {\end{lrbox}\noindent{\centering
    \IfDecimal{\ctabargstr}%
      {\scalebox{\ctabarg}{\usebox{\ctabsavebox}}}%
      {\usebox{\ctabsavebox}}\par}%
   \vspace{2mm}}

% Emphasis: bold rather than italic (italic is hard to read on screen).
\renewcommand{\emph}[1]{\textbf{#1}}

% Proof heading: bold (override amsthm's default \itshape).
\makeatletter
\renewenvironment{proof}[1][\proofname]{\par
  \pushQED{\qed}%
  \normalfont \topsep6\p@\@plus6\p@\relax
  \trivlist
  \item[\hskip\labelsep
        \bfseries
    #1\@addpunct{.}]\ignorespaces
}{%
  \popQED\endtrivlist\@endpefalse
}
\makeatother

% Three signature pedagogical boxes:
%   motivation (blue)  : opening hooks, contextual framing
%   derivation (green) : supplementary derivations beyond the main flow
%   suppremark (orange): instructor commentary, study tips, broader perspective
\newtcolorbox{motivation}[1][Motivation]{breakable, enhanced, colback=blue!5,
  colframe=blue!50!black, coltitle=white,
  title=\textbf{#1}, fonttitle=\bfseries,
  attach boxed title to top left={yshift=-2mm,xshift=3mm},
  boxed title style={colback=blue!50!black}}
\newtcolorbox{derivation}[1][Derivation]{breakable, enhanced, colback=green!5,
  colframe=green!40!black, coltitle=white,
  title=\textbf{#1}, fonttitle=\bfseries,
  attach boxed title to top left={yshift=-2mm,xshift=3mm},
  boxed title style={colback=green!40!black}}
\newtcolorbox{suppremark}[1][Remark]{breakable, enhanced, colback=orange!5,
  colframe=orange!60!black, coltitle=white,
  title=\textbf{#1}, fonttitle=\bfseries,
  attach boxed title to top left={yshift=-2mm,xshift=3mm},
  boxed title style={colback=orange!60!black}}

% Custom math operators used throughout the course.
\usepackage{dsfont}
\DeclareMathOperator\prb{\mathsf{P}}
\DeclareMathOperator\expc{\mathsf{E}}
\DeclareMathOperator\var{var}
\DeclareMathOperator\cov{cov}
\DeclareMathOperator\cor{cor}
\DeclareMathOperator\indc{\mathds{1}}
\DeclareMathOperator\dom{dom}
\DeclareMathOperator\epi{epi}
\DeclareMathOperator\diag{diag}
\DeclareMathOperator\rk{rank}
\DeclareMathOperator\nul{null}
\DeclareMathOperator\col{col}
\DeclareMathOperator\img{img}
\let\ker\relax
\DeclareMathOperator\ker{ker}
\DeclareMathOperator\tr{tr}
\DeclareMathOperator\df{df}

\newcommand{\dd}{\mathrm{d}}

\usepackage{parskip}
\setlength{\parskip}{3pt}\setlength{\jot}{4pt}

\usepackage{titlesec}
\titleformat{\section}{\color{hdrBG}\normalfont\Large\bfseries}{\thesection}{0.7em}{}
\titleformat{\subsection}{\color{accent}\normalfont\large\bfseries}{\thesubsection}{0.6em}{}
\titlespacing*{\section}{0pt}{4pt}{2pt}
\titlespacing*{\subsection}{0pt}{4pt}{1pt}
